
\documentclass[12pt,a4paper]{article}

\usepackage[margin=1in]{geometry}
\usepackage{amsmath,amssymb}
\usepackage{booktabs}
\usepackage{array}
\usepackage{listings}
\usepackage{float}
\usepackage[T1]{fontenc}
\usepackage{lmodern}
\usepackage{hyperref}

\hypersetup{
    colorlinks=true,
    linkcolor=black,
    urlcolor=blue,
    citecolor=black
}

\lstset{
    language=Python,
    basicstyle=\ttfamily\small,
    breaklines=true,
    frame=single,
    columns=fullflexible,
    showstringspaces=false
}

\title{\textbf{Chapter 5: Resampling Methods}\\
Exercises 1 and 9}
\author{Student Name: \underline{\hspace{6cm}}\\
Student ID: \underline{\hspace{6cm}}}
\date{\today}

\begin{document}

\maketitle

\section*{Introduction}

This report presents solutions to Exercises 1 and 9 from Chapter 5,
\textit{Resampling Methods}, of \textit{An Introduction to Statistical
Learning with Applications in Python (ISLP)}. Exercise 1 is a conceptual
derivation involving the minimum-variance combination of two random
variables. Exercise 9 applies bootstrap methods to the Boston housing
data set to estimate the mean, median, and tenth percentile of the
\texttt{medv} variable.

The calculations in Exercise 9 were implemented in Python using the
\texttt{ISLP} package. The bootstrap calculations use 10,000 bootstrap
samples, with the random number generator initialized using seed 1.

\section{Exercise 1}

\subsection*{Question}

Using the basic statistical properties of the variance, as well as
single-variable calculus, derive equation (5.6). In other words, show
that the value of $\alpha$ given by equation (5.6) minimizes
\[
\operatorname{Var}\left(\alpha X+(1-\alpha)Y\right).
\]

\subsection*{Solution}

Let
\[
V(\alpha)=\operatorname{Var}\left(\alpha X+(1-\alpha)Y\right).
\]

Define
\[
\sigma_X^2=\operatorname{Var}(X), \qquad
\sigma_Y^2=\operatorname{Var}(Y),
\]
and
\[
\sigma_{XY}=\operatorname{Cov}(X,Y).
\]

Using the variance formula
\[
\operatorname{Var}(aX+bY)
=
a^2\operatorname{Var}(X)
+b^2\operatorname{Var}(Y)
+2ab\operatorname{Cov}(X,Y),
\]
we obtain
\[
V(\alpha)
=
\alpha^2\sigma_X^2
+(1-\alpha)^2\sigma_Y^2
+2\alpha(1-\alpha)\sigma_{XY}.
\]

Expanding the terms gives
\[
V(\alpha)
=
\alpha^2\sigma_X^2
+
(1-2\alpha+\alpha^2)\sigma_Y^2
+
(2\alpha-2\alpha^2)\sigma_{XY}.
\]

Collecting terms according to powers of $\alpha$,
\[
V(\alpha)
=
\alpha^2
(\sigma_X^2+\sigma_Y^2-2\sigma_{XY})
+
\alpha(-2\sigma_Y^2+2\sigma_{XY})
+
\sigma_Y^2.
\]

To find the value of $\alpha$ that minimizes the variance, differentiate
with respect to $\alpha$:
\[
\frac{dV}{d\alpha}
=
2\alpha(\sigma_X^2+\sigma_Y^2-2\sigma_{XY})
-2\sigma_Y^2+2\sigma_{XY}.
\]

At the minimum,
\[
\frac{dV}{d\alpha}=0.
\]

Therefore,
\[
2\alpha(\sigma_X^2+\sigma_Y^2-2\sigma_{XY})
-2\sigma_Y^2+2\sigma_{XY}=0.
\]

Dividing by 2,
\[
\alpha(\sigma_X^2+\sigma_Y^2-2\sigma_{XY})
=
\sigma_Y^2-\sigma_{XY}.
\]

Hence,
\[
\boxed{
\hat{\alpha}
=
\frac{\sigma_Y^2-\sigma_{XY}}
{\sigma_X^2+\sigma_Y^2-2\sigma_{XY}}
}
\]

which is the required result.

\subsection*{Verification that the solution is a minimum}

We take the second derivative:
\[
\frac{d^2V}{d\alpha^2}
=
2(\sigma_X^2+\sigma_Y^2-2\sigma_{XY}).
\]

Notice that
\[
\sigma_X^2+\sigma_Y^2-2\sigma_{XY}
=
\operatorname{Var}(X-Y).
\]

Since variance is non-negative,
\[
\operatorname{Var}(X-Y)\geq 0.
\]

Therefore,
\[
\frac{d^2V}{d\alpha^2}\geq 0.
\]

Thus, the value of $\hat{\alpha}$ obtained above minimizes
$\operatorname{Var}(\alpha X+(1-\alpha)Y)$.

\section{Exercise 9: Boston Housing Data}

\subsection*{Data and Method}

The Boston housing data set contains 506 observations. The variable
\texttt{medv} represents the median value of owner-occupied homes for
Boston census tracts, measured in thousands of dollars.

For parts (c), (f), and (h), bootstrap sampling was used. Each bootstrap
sample contains 506 observations selected with replacement from the
original data. The relevant statistic was calculated for each of 10,000
bootstrap samples. The bootstrap standard error was then estimated as
the standard deviation of the resulting bootstrap statistics.

\subsection*{Python setup}

The data were loaded as follows:

\begin{lstlisting}
from ISLP import load_data
import numpy as np

Boston = load_data("Boston")
Boston.head()
\end{lstlisting}

\subsection*{(a) Estimate the population mean}

The population mean is estimated using the sample mean:
\[
\hat{\mu}=\bar{x}.
\]

In Python:

\begin{lstlisting}
mu_hat = Boston['medv'].mean()
print(mu_hat)
\end{lstlisting}

The result is
\[
\boxed{\hat{\mu}=22.5328063}.
\]

Therefore, the estimated population mean of \texttt{medv} is
approximately $22.53$ thousand dollars.

\subsection*{(b) Estimate the standard error of $\hat{\mu}$}

The standard error of the sample mean is estimated using
\[
SE(\hat{\mu})=\frac{s}{\sqrt{n}},
\]
where $s$ is the sample standard deviation and $n$ is the number of
observations.

The Python implementation is:

\begin{lstlisting}
se_mu = Boston['medv'].std() / np.sqrt(Boston.shape[0])
print(se_mu)
\end{lstlisting}

The result is
\[
\boxed{SE(\hat{\mu})=0.4088611}.
\]

This means that the sampling variability of the estimated mean is about
$0.409$ thousand dollars.

\subsection*{(c) Bootstrap estimate of the standard error}

The bootstrap procedure was used to estimate the standard error of the
mean. The procedure repeatedly draws samples of size 506 with
replacement and calculates the mean of each sample.

\begin{lstlisting}
rng = np.random.default_rng(1)

bootstrap_means = []

for i in range(10000):
    sample = rng.choice(
        Boston['medv'],
        size=len(Boston),
        replace=True
    )

    bootstrap_means.append(sample.mean())

bootstrap_se = np.std(bootstrap_means, ddof=1)

print(bootstrap_se)
\end{lstlisting}

The bootstrap standard error is
\[
\boxed{SE_{\text{bootstrap}}(\hat{\mu})=0.4097871}.
\]

The analytical standard error from part (b) was $0.4088611$, while the
bootstrap estimate is $0.4097871$. The two estimates are extremely close.
This indicates that the bootstrap provides a good estimate of the
sampling variability of the sample mean in this data set.

\subsection*{(d) 95\% confidence interval for the population mean}

The exercise suggests the approximate confidence interval
\[
\left[
\hat{\mu}-2SE(\hat{\mu}),
\hat{\mu}+2SE(\hat{\mu})
\right].
\]

Using the bootstrap standard error from part (c),
\[
\hat{\mu}=22.5328063
\]
and
\[
SE_{\text{bootstrap}}=0.4097871.
\]

Therefore,
\[
22.5328063-2(0.4097871)=21.7132322
\]
and
\[
22.5328063+2(0.4097871)=23.3523805.
\]

Thus, the approximate 95\% confidence interval is
\[
\boxed{[21.7132,\;23.3524]}.
\]

The Python calculation is:

\begin{lstlisting}
lower = mu_hat - 2 * bootstrap_se
upper = mu_hat + 2 * bootstrap_se

print(lower, upper)
\end{lstlisting}

Hence, we estimate that the population mean of \texttt{medv} lies
approximately between $21.71$ and $23.35$ thousand dollars.

\subsection*{(e) Estimate the population median}

The population median is estimated using the sample median:

\begin{lstlisting}
mu_med = Boston['medv'].median()
print(mu_med)
\end{lstlisting}

The result is
\[
\boxed{\hat{\mu}_{med}=21.2}.
\]

Therefore, the estimated population median of \texttt{medv} is
$21.2$ thousand dollars.

\subsection*{(f) Bootstrap standard error of the median}

Unlike the sample mean, there is no simple formula of the same form for
the standard error of the median. Therefore, the bootstrap is used.

\begin{lstlisting}
bootstrap_medians = []

for i in range(10000):
    sample = rng.choice(
        Boston['medv'],
        size=len(Boston),
        replace=True
    )

    bootstrap_medians.append(np.median(sample))

bootstrap_median_se = np.std(
    bootstrap_medians,
    ddof=1
)

print(bootstrap_median_se)
\end{lstlisting}

The bootstrap estimate of the standard error is
\[
\boxed{SE_{\text{bootstrap}}(\hat{\mu}_{med})=0.3813021}.
\]

This shows that the estimated median has sampling variability of about
$0.381$ thousand dollars. The bootstrap is useful here because it
provides a practical estimate of uncertainty without requiring a simple
analytical formula for the median's standard error.

\subsection*{(g) Estimate the tenth percentile}

The tenth percentile is estimated using NumPy's
\texttt{percentile()} function:

\begin{lstlisting}
np.percentile(Boston['medv'], 10)
\end{lstlisting}

The result is
\[
\boxed{\hat{\mu}_{0.1}=12.75}.
\]

Thus, approximately 10\% of the observations have a \texttt{medv} value
below $12.75$ thousand dollars.

\subsection*{(h) Bootstrap standard error of the tenth percentile}

The bootstrap was again used, this time calculating the tenth percentile
for every bootstrap sample.

\begin{lstlisting}
bootstrap_percentiles = []

for i in range(10000):
    sample = rng.choice(
        Boston['medv'],
        size=len(Boston),
        replace=True
    )

    bootstrap_percentiles.append(
        np.percentile(sample, 10)
    )

bootstrap_percentile_se = np.std(
    bootstrap_percentiles,
    ddof=1
)

print(bootstrap_percentile_se)
\end{lstlisting}

The resulting bootstrap standard error is
\[
\boxed{SE_{\text{bootstrap}}(\hat{\mu}_{0.1})=0.4999335}.
\]

This indicates that the estimated tenth percentile has a sampling
variability of approximately $0.500$ thousand dollars. The standard
error is larger than the standard error of the mean because percentile
estimation, especially toward the lower tail of a distribution, can be
more variable.

\section{Summary of Results}

\begin{table}[H]
\centering
\begin{tabular}{>{\raggedright\arraybackslash}p{1.2cm}
                >{\raggedright\arraybackslash}p{6.0cm}
                >{\raggedleft\arraybackslash}p{3.0cm}}
\toprule
\textbf{Part} & \textbf{Quantity} & \textbf{Result} \\
\midrule
9(a) & Sample estimate of population mean & 22.5328 \\
9(b) & Analytical standard error of mean & 0.4089 \\
9(c) & Bootstrap standard error of mean & 0.4098 \\
9(d) & Approximate 95\% confidence interval & [21.7132, 23.3524] \\
9(e) & Sample estimate of population median & 21.2 \\
9(f) & Bootstrap standard error of median & 0.3813 \\
9(g) & Tenth percentile & 12.75 \\
9(h) & Bootstrap standard error of tenth percentile & 0.4999 \\
\bottomrule
\end{tabular}
\caption{Results obtained from the Boston housing data.}
\end{table}

\section{Conclusion}

Exercise 1 demonstrates how basic variance properties and differentiation
can be used to obtain the minimum-variance weight for a combination of
two random variables. The resulting estimator is
\[
\hat{\alpha}
=
\frac{\sigma_Y^2-\sigma_{XY}}
{\sigma_X^2+\sigma_Y^2-2\sigma_{XY}}.
\]

Exercise 9 demonstrates the practical value of the bootstrap. For the
Boston housing data, the bootstrap estimate of the standard error of the
mean ($0.4098$) was very close to the analytical standard error
($0.4089$). The bootstrap was also used to estimate the standard errors
of the median and tenth percentile, for which a simple standard-error
formula is not directly available in the exercise.

Overall, the results show how resampling can be used to quantify the
uncertainty associated with statistical estimates using repeated samples
drawn from the observed data.

\section*{Reference}

James, G., Witten, D., Hastie, T., Tibshirani, R., \& Taylor, J.
\textit{An Introduction to Statistical Learning with Applications in
Python}. Springer, 2023.

\end{document}
